\documentclass{article}
\usepackage[utf8]{inputenc}
\usepackage[T1]{fontenc}
\usepackage[margin=2cm]{geometry}
\usepackage{hyperref}
\sloppy
\begin{document}
\section*{Historical source}
Encyclopaedia Britannica scan, source row 1250. Parse 3.5 extracted the text. Super 3.5 VL corrected it from an enlarged crop of each block. It can still contain recognition errors.
Page headers, footers, and page numbers are removed.

\subsection*{Untitled continuation (E0099)}
\textit{Continues an entry from an earlier page.}

Coin,
Coinage.
According to the decuple pro-
portion, mentioned by Livy and
Julius Pollux, worth

According to the proportion mentioned by Tacitus, and which afterwards obtained, whereby the aureus exchanged for 25 denarii, its value

Coin, in architecture, a kind of dye cut diagonally, after the manner of a flight of a stair-case, serving at bottom to support columns in a level, and at top to correct the inclination of an entablature supporting a vault.

COIN is also used for a solid angle composed of two surfaces inclined towards each other, whether that angle be exterior, as the coin of a wall, a tree, \&c. or interior, as the coin of a chamber or chimney. See Quoin.

\subsection*{COINAGE (E0100)}
COINAGE, or Coining, the art of making money, as performed either by the hammer or mill.

Formerly the fabric of coins was different from what it is at present. They cut a large plate of metal into several little squares, the corners of which were cut off with shears. After having shaped these pieces, so as to render them perfectly conformable, in point of weight, to the standard piece, they took each piece in hand again, to make it exactly round by a gentle hammering. This was called a planchet, and was fit for immediate coining. Then engravers prepared, as they still do, a couple of steel mauls in form of dies, cut and terminated by a flat surface, rounded off at the edges. They engraved or stamped on it the hollow of a head, a cross, a fleur-de-lis, or any other figure, according to the custom of the times, with a short legend. As one of these dies was to remain dormant, and the other movable, the former ended in a square prism, that it might be introduced into the square hole of the block, which, being fixed very fast, kept the dye as steady as any vice could have done. The planchet of metal was horizontally laid upon this inferior maul, to receive the stamp of it on one side, and that of the upper dye, wherewith it was covered, on the other. This movable dye, having its round engraved surface resting upon the planchet, had at its opposite extremity a flat square, and larger surface, upon which they gave several heavy blows, with a hammer of an enormous size, till the double stamp was sufficiently, in relief, impressed on each side of the planchet. This being finished, was immediately succeeded by another, and they thus became a standard coin, which had the degree of fineness of the weight and mark determined by the judgment of the inspectors, to make it good current money. The strong tempering which was and is still given to the two dyes, rendering them capable of bearing those repeated blows. Coining has been considerably improved and rendered expeditious, by several ingenious machines, and by a wide application of the useful physical experiments to the methods of fixing, drying, and stamping the different metals.

The three finest instruments the mint-man ufes, are the laminating engine; the machine for making the impreffions on the edges of coins ; and the mill.

After they have taken the laminae, or plates of me-

tal, out of the mould into which they are cast, they do not beat them on the anvil, as was formerly done, but they make them pas and repass between the several rollers of the laminating engine, which being gradually brought closer and closer to each other, frequently give the lamina its uniform and exact thickness. Instead of dividing the lamina into small squares, they at once cut clean out of it as many planchets as it can contain, by means of a sharp steel trepan, of a roundish figure, hollow within, and of a proportionable diameter, to shape and cut off the piece at one and the same time. After these planchets have been prepared and weighed with standard pieces, filed or forged to get off the superfluous part of the metal, and then boiled and made clean, they arrive, at last, at the machine (fig. 1.), which marks them upon the edge; and finally, the mill (fig. 2.), which, squeezing each of them finely between the two dyes, brought near each other with one blow, forces the two surfaces or fields of the piece to fill exactly all the vacancies of the two figures engraved hollow. The engine which serves to laminate lead, gives a sufficient notion of that which serves to flatten gold and silver. laminae between rollers of a lesser size.

The principal pieces of the machine (fig. 1.), to stamp coins on the edge, are two steel laminae, about a line thick. One half of the legend, or of the ring, is engraved on the thickness of one of the laminae, and the other half on the thickness of the other; and these two laminae are straight, although the planchet marked with them be circular.

When they stamp a planchet, they first put it between the laminæ in such a manner, as that the being each of them laid flat upon a copper-plate, which is fastened upon a very thick wooden table, and the planchet being likewise laid flat upon the same plate, the edge of the planchet may touch the two laminæ on each side, and in their thick part.

One of these laminae is immoveable, and fastened with several screws; the other slides by means of a dented wheel, which takes into the teeth that are on the surface of the lamina. This sliding lamina makes the planchet turn in such a manner, that it remains stamped on the edge, when it has made one turn. Only crown and half-crown pieces can bear the impression of letters on the thickness of their edges.

The coining engine or mill is so handy (fig. 2.), that a single man may stamp twenty thousand planchets in one day: gold, silver, and copper planchets, are all of them coined with a mill, to which the coining squares (fig. 3.), commonly called dyes, are fastened; that of the face under, in a square box garnished with male and female screws, to fix and keep it steady; and the other above, in a little box garnished with the same screws, to fasten the coining square. The planchet is laid flat on the square of the effigy, which is dormant; and they immediately pull the bar of the mill by its cords, which causes the screw set within it to turn. This enters into the female screw, which is in the body of the mill, and turns with so much strength, that by pulling the upper square upon that of the effigy, the planchet, violently pressed between both squares, receives the impression of both at one pull, and in the twinkling of an eye.


\section*{Source record}
\begin{description}
\item[Scan ID] \texttt{britannica:309f765b1ecc83e0a7f7762e44b4a927a7555443:train:1250}
\item[Image file] \texttt{page-1250-2e0c75f85418.jpg}
\item[Image SHA-256] \texttt{2e0c75f854185185ace0f165eea2b904222208282d93a1c66b97c9b5c42f5c95}
\item[Dataset] NationalLibraryOfScotland/encyclopaedia\_britannica\_illustrated at revision \texttt{309f765b1ecc83e0a7f7762e44b4a927a7555443} (CC0-1.0)
\item[Extraction] model \texttt{nvidia/nemotron-3.5-parse}, raw reply \texttt{extract/raw/page-1250-2e0c75f85418.json}
\item[Raw reply SHA-256] \texttt{4c755df1a9b6eca65fcc1837742bbee91e53b6ff4b89e84fc78934c1ee517772}
\item[Correction] model \texttt{nvidia/Nemotron-3.5-Super-VL}, block crops, raw replies \texttt{correct/raw/crop\_nothink/page-1250-2e0c75f85418-b*.json}; super\_vl: 13
\item[Running header] COI I 130 | COI
\end{description}
\end{document}
